\chapter{인플레이션, 구조 형성, 초기 조건의 소실}

\section{오메가와 평탄성의 증거}

큰 \(\Omega\)에 관한 한 학생의 질문으로 논의는 프리드만-르메트르-로버트슨-워커 방정식으로 되돌아간다. 복사 지배기가 한참 지난 현재 우주에서는 팽창 법칙을 진공 에너지 항과 물질 항으로 나타내며, 복사는 이제 매우 작으므로 생략한다.
\begin{equation}
H^2=\left(\frac{\dot a}{a}\right)^2=\frac{8\pi G}{3}\left(\rho_0+\rho_m\right)-\frac{k}{a^2}.
\end{equation}
진공 항은 희석되지 않는다. 물질 항은 희석된다.
\begin{equation}
\rho_m=\frac{M}{a^3}.
\end{equation}
\(\Omega\)는 팽창 방정식에서 각 성분이 차지하는 비율이다.

\(H^2\)로 나누면 정규화된 항들의 합은 1이 된다. 정규화된 표기법으로 쓰면
\begin{equation}
\begin{aligned}
\Omega_0&=\frac{8\pi G\rho_0}{3H^2},\\
\Omega_m&=\frac{8\pi G\rho_m}{3H^2},\\
\Omega_k&=-\frac{k}{a^2H^2},
\end{aligned}
\end{equation}
따라서
\begin{equation}
\Omega_0+\Omega_m+\Omega_k=1.
\end{equation}
\editorialnote{이 대목의 음성과 칠판 내용은 선명하지 않고, 강연자는 진공 기여분을 가리키는 명칭을 몇 차례 바꾸어 쓴다. 여기서는 \(H^2\)로 나눈 뒤의 논지를 정규화된 표기법으로 정리했다.}
공간이 평탄하면 \(k=0\)이므로
\begin{equation}
\Omega_0+\Omega_m=1.
\end{equation}
물질과 진공의 합이 1보다 크면 곡률은 양수이고, 1보다 작으면 곡률은 음수이다.

그러면 관측 논리는 간단하다. 보통 물질과 암흑 물질을 측정한다. 허블 팽창을 측정한다. 매우 큰 규모에서 기하학적으로 곡률을 측정한다. 또한 우주가 가속 팽창하고 있음에 주목한다. 진공 항의 부호 자체가 \(H\)인 것이 핵심은 아니다. 핵심은 \(\dot a/a\)가 상수에 가까워지는 경향이며, 이는 스케일 인자가 지수 팽창 가지에 접근한다는 뜻이다.
\begin{equation}
a(t)\propto e^{Ht}.
\end{equation}
여러 측정 결과는 서로 교차 검증하며, 우주가 몇 퍼센트 이내의 오차로 거의 평탄하다는 것을 보여 준다.

\begin{classroomqa}
\lecturetimestamp{00:10:50}
\audiencequestion{물질에 암흑 물질도 포함되나요?}
\lecturerresponse{그렇다. 여기서 논의하는 물질 항에는 암흑 물질이 포함된다. 복사도 더할 수 있지만, 오늘날 복사가 차지하는 비율은 무시할 수 있을 만큼 작다.}
\end{classroomqa}

\begin{classroomqa}
\lecturetimestamp{00:11:11}
\audiencequestion{허블 상수 말고 암흑 에너지를 측정할 다른 방법이 있나요?}
\lecturerresponse{핵심은 \(H\) 자체의 값이 아니라 팽창이 지수 법칙에 가까워지는 양상이다. 이를 물질 측정값 및 평탄성의 기하학적 측정 결과와 함께 이용하면 진공 항을 추론할 수 있다.}
\end{classroomqa}

\section{평탄성에서 구조 형성으로}

우주는 왜 이렇게 평탄할까? 인플레이션이 그 해답이라는 제안이 있다. 평탄하다는 것은 수십억, 실은 수백억 광년 규모에서 우주가 매우 크고, 매우 매끄러우며, 매우 균질하다는 뜻이다. 지수 팽창은 한 영역을 늘려 그렇게 보이게 할 수 있다.

하지만 좋은 일도 지나치면 문제가 된다. 완벽하게 매끄러운 우주는 계속 완벽하게 매끄럽다. 중력은 완벽한 대칭성에서 덩어리를 만들어 내지 않는다. 지금 구조가 존재한다면, 그 구조는 과거의 작은 비균질성에서 시작했어야 한다. 현재 우주를 거꾸로 추적해 보면, 투명해졌을 당시의 비균질성은 대략 \(10^5\)분의 1이다.
\begin{equation}
\frac{\delta\rho}{\rho}\sim 10^{-5}.
\end{equation}
우주는 극도로 매끄러웠지만 완벽하게 균질하지는 않았다.

그 뒤에는 중력이 공기로 가득 찬 방 안의 보통 힘과는 다른 방식으로 작용한다. 과밀 영역은 물질을 자기 쪽으로 끌어당긴다. 주변 영역은 물질이 줄어들므로 저밀도 영역은 더 저밀해지고 과밀 영역은 더 과밀해진다. 중력은 대비를 없애는 대신 강화한다.

자연스럽게 드는 질문은 왜 이런 일이 더 일찍 일어나지 않았느냐는 것이다. 답은 압력이다. 더 이른 시기에는 복사압이 컸다. 다음 식을 쓰면
\begin{equation}
p=w\rho,
\end{equation}
익숙한 경우들은 다음과 같다.
\begin{equation}
w_{\rm rad}=\frac13,\qquad
w_m=0,\qquad
w_{\rm vac}=-1.
\end{equation}
양의 압력은 균질화한다. 요동을 밀어내 매질을 매끄럽게 만든다. 큰 규모에서 압력이 없는 것으로 다루는 물질에는 그런 지지력이 없다. 압력이 있으면 중력의 작용이 억제된다. 압력이 무시할 수 있을 만큼 작아지면 중력이 우세해진다.

암흑 물질은 사건의 순서에서 핵심적인 역할을 한다. 처음 붕괴하는 덩어리는 암흑 물질 헤일로다. 그다음 보통의 바리온 물질이 헤일로 안으로 떨어지고, 바리온은 에너지를 소산할 수 있으므로 중심에 더 가까이 모인다. 암흑 물질은 같은 방식으로 에너지를 소산하지 않아서 헤일로 전반에 퍼진 채 남는다. 이 수축의 시작은 투명해지기보다 약간 이른데, 시간으로 대략 10배 빠르다. 그래서 우주 마이크로파 배경에는 그 시작이 직접 나타나지 않는다.

짧게 덧붙이면 이후의 화학적 역사가 정리된다. 헬륨보다 무거운 원소는 거의 모두 초기 우주론이 아니라 항성 진화의 산물이다. 초기 우주에서는 사실상 수소와 헬륨, 그리고 몇 가지 가벼운 원자핵이 소량 만들어진다.

하늘에서 관측하는 것은 질량 밀도 자체가 아니라 우주 마이크로파 배경의 온도 비등방성이다. 밀도 요동이 여기에 기여한다. 속도 요동도 기여한다. 과밀 영역은 더 깊은 중력 퍼텐셜 우물에서 빠져나오는 광자가 에너지를 더 잃기 때문에 약간 더 차갑게 보이는 경향이 있으며, 여기에 기체의 운동에 따른 도플러 편이가 더해진다. 전체 수치 분석은 복잡하다. 확실하게 말할 수 있는 내용은 더 간단하다. 우주가 투명해졌을 무렵 에너지 밀도는 \(10^5\)분의 1 정도 달라졌다.

\begin{classroomqa}
\lecturetimestamp{00:20:29}
\audiencequestion{중력은 왜 덩어리를 더 일찍 키우지 않았나요?}
\lecturerresponse{그보다 이른 시기에는 우주가 복사 지배 상태였고, 복사는 매질을 균질화할 만큼 충분한 양의 압력을 지니고 있었다. 복사가 물질에 비해 중요하지 않게 된 뒤에야 중력이 덩어리의 성장을 지배하기 시작했다.}
\end{classroomqa}

\begin{classroomqa}
\lecturetimestamp{00:24:35}
\audiencequestion{붕괴하려면 소산 메커니즘이 필요하지 않나요?}
\lecturerresponse{보통의 바리온 물질이 중심을 향해 가라앉으려면 소산이 필요하다. 그러나 암흑 물질은 그런 종류의 충돌 소산 없이도 헤일로를 형성하기 시작할 수 있다. 팽창 자체와 관련된 우주론적 감쇠도 있으며, 이것은 나중에 중요해진다.}
\end{classroomqa}

\begin{classroomqa}
\lecturetimestamp{00:30:34}
\audiencequestion{속도 요동도 있었나요? 그래야 나중에 은하가 각운동량을 얻을 수 있지 않나요?}
\lecturerresponse{그렇다. 원시적인 요동 양상에는 밀도 요동과 속도 요동이 모두 있었다. 하늘에서 관측되는 온도 분포는 밀도 변화와 기체의 운동이 함께 만들어 낸 것이다.}
\end{classroomqa}

\section{인플라톤과 평탄부}

이제 소개할 메커니즘에는 스칼라장 \(\phi\)가 등장한다. 이 장의 에너지 밀도는 세 부분으로 이루어진다.
\begin{equation}
\rho_\phi=\frac12\dot\phi^2+\frac12(\nabla\phi)^2+V(\phi).
\end{equation}
시간에 따른 변화가 기여한다. 공간에 따른 변화도 기여한다. 스칼라장은 진정한 퍼텐셜 에너지 \(V(\phi)\)도 지닐 수 있다. 이 부분은 이국적인 개념이 아니다. 익숙한 예로 힉스장이 있다. 여기서 중요한 것은 퍼텐셜의 형태다.

퍼텐셜은 어느 범위에서 상당히 평탄하면서 약간 아래로 기울어져 있다가, 임계값에서 바닥이 급격히 내려가야 한다. 정확히 0까지 떨어지는 것은 아니고, 매우 작은 값까지 내려간다. 플랑크 단위로 현재의 진공 에너지는 대략 다음 크기다.
\begin{equation}
10^{-123}.
\end{equation}
평탄부에서는 이보다 훨씬 컸지만, 그래도 플랑크 규모보다는 훨씬 작았다.

\begin{figure}[t]
\centering
\includegraphics[width=0.78\textwidth]{lecture_08_figure_03.png}
\caption{넓은 평탄부와 급경사 하강부, 낮은 최솟값을 갖는 인플라톤 퍼텐셜의 정성적 형태.}
\lectureframe{00:43:23}
\end{figure}

이 그래프는 공간상의 위치를 나타내지 않는다. 장의 값을 나타낸다. 우주가 평탄부의 높은 곳에서 시작했다고 말하는 것은 매우 넓은 영역에서 장의 값이 그 높은 곳에 있었다는 뜻일 뿐이다. 처음에 \(\phi\)가 장소마다 조금 달랐는지는 중요하지 않다. 인플레이션이 그 차이를 엄청나게 늘리기 때문이다. 충분히 팽창한 뒤에는 어느 국소 영역도 거의 균일해 보인다. \editorialnote{칠판 화면에서 축의 표기는 읽을 수 없지만, 강연 내용을 보면 세로축은 퍼텐셜 에너지 \(V(\phi)\)이고 가로축은 장의 값 \(\phi\)임을 알 수 있다.}

장이 넓고 완만하게 기울어진 퍼텐셜 부분에 머무는 동안 퍼텐셜 에너지는 진공 에너지처럼 작용한다. \(\phi\)가 완전히 고정되어 있다면 진공 에너지와 정확히 똑같이 작용할 것이다. 생략된 수치 계수는 제쳐 두면, 강연에서 제시한 관계는 간단히 다음과 같다.
\begin{equation}
H\propto \sqrt{V(\phi)},
\end{equation}
그러면 거의 지수적인 법칙이 나온다.
\begin{equation}
a(t)\propto e^{Ht}.
\end{equation}
현재 관측할 수 있는 거의 모든 현상은 장이 가장자리를 넘어간 뒤에 일어난다. 그 하강 과정에서 평탄부의 큰 진공 에너지와 비슷한 에너지가 복사, 보통 물질, 암흑 물질, 그리고 오늘날 남아 있는 미소한 진공 에너지로 전환된다.

\begin{classroomqa}
\lecturetimestamp{00:44:07}
\audiencequestion{\(\phi\)는 무엇 위에 정의된 장이며, 아인슈타인 방정식은 여전히 쓰이나요?}
\lecturerresponse{그렇다. \(\phi\)는 시공간 위에 정의된 장이며, 같은 프리드만 방정식을 계속 사용한다. 이 영역에서는 장이 충분히 천천히 변한다면 우변의 에너지 밀도는 주로 \(V(\phi)\)가 제공한다.}
\end{classroomqa}

\begin{classroomqa}
\lecturetimestamp{00:47:04}
\audiencequestion{진공 에너지처럼 작용하는 인플라톤 에너지는 어떻게 물질과 복사로 바뀌나요?}
\lecturerresponse{장이 바닥까지 굴러 내려가 진동하면, 그 진동은 장의 양자에 해당한다. 그 양자가 불안정하면 광자, 전자, 양전자 및 다른 익숙한 입자로 붕괴할 수 있다.}
\end{classroomqa}

\section{느린 굴림과 허블 마찰}

장이 가장자리를 넘어간 뒤 왜 그저 영원히 출렁이지 않는가? 왜 앞뒤로 진동하다가 심지어 다시 올라가지 않는가? 여기서 제시하는 답은 우주론적 감쇠 메커니즘이다. 이것은 미시적인 의미의 실제 마찰은 아니지만, 운동 방정식에서는 마찰처럼 작용한다.

공간 기울기는 무시하고 장을 균질하다고 보자. 그러면 운동 항과 퍼텐셜 항은 \(T-V\) 역학을 따르는 하나의 자유도로 보이지만, 에너지 밀도이므로 물리적 부피 \(a^3\)를 곱해야 한다. 따라서 균질한 유효 라그랑지언은 다음 형태가 된다.
\begin{equation}
L_{\mathrm{hom}}=a^3(t)\left[\frac12\dot\phi^2-V(\phi)\right].
\end{equation}
\editorialnote{칠판의 맨 위 줄은 일부가 잘려 있다. 아래에 적힌 확실한 운동 방정식을 재현하는 데 필요한 범위에서만 이 균질 라그랑지언을 복원했다.}

칠판에서 확실하게 확인되는 방정식은 다음과 같다.
\begin{equation}
\frac{d}{dt}\left(a^3\dot\phi\right)=-a^3\frac{\partial V}{\partial\phi}.
\end{equation}
시간 미분을 전개하면
\begin{equation}
a^3\ddot\phi+3a^2\dot a\,\dot\phi=-a^3\frac{\partial V}{\partial\phi},
\end{equation}
그리고 \(a^3\)로 나누면 다음을 얻는다.
\begin{equation}
\ddot\phi+3H\dot\phi=-\frac{\partial V}{\partial\phi}.
\end{equation}
구어로 설명하는 과정에서 음의 부호를 명시적으로 바로잡는다. 힘은 퍼텐셜 에너지의 미분값에 음수를 붙인 것이다.

\begin{figure}[t]
\centering
\includegraphics[width=0.78\textwidth]{lecture_08_figure_05.png}
\caption{허블 마찰 항 \(3H\dot\phi\)이 나타난 균질 인플라톤 운동 방정식.}
\lectureframe{00:57:10}
\end{figure}

\(3H\dot\phi\)라고 부르는 항이 바로 허블 마찰이다. 이는 전적으로 \(a^3\) 인자를 미분할 때 생긴다. 그러나 동역학적으로 보면 물속을 떨어지는 돌에 작용하는 항력과 비슷하다. 힘 항과 가속도 항, 속도에 선형으로 비례하는 항이 있다. 장은 빠르게 종단 속도에 도달한 뒤 평탄부를 따라 천천히 움직인다. 이것이 느린 굴림 영역이다. 기울기가 완만하고 장이 천천히 움직이는 동안에는 퍼텐셜 에너지가 천천히 변하므로 팽창은 지수 팽창에 가깝게 유지된다.

최솟값 근처에서 장은 진동한다. 그 진동은 인플라톤 양자에 해당한다. 그 양자가 불안정하면 붕괴하며, 이 과정에서 평탄부의 에너지가 후기 우주의 뜨거운 물질과 복사 성분으로 전환된다.

\begin{classroomqa}
\lecturetimestamp{00:50:48}
\audiencequestion{장은 왜 영원히 진동하지 않으며, 감쇠는 무엇이 일으키나요?}
\lecturerresponse{팽창 자체가 마찰과 비슷한 항을 만든다. 균질 라그랑지언의 \(a^3\) 인자가 \(3H\dot\phi\) 항을 만들어 내며, 이 항은 점성 항력처럼 작용한다.}
\end{classroomqa}

\section{양자 요동과 지평선 이탈}

이제 겉보기에는 문제가 생긴다. 인플레이션이 우주를 극도로 매끄럽게 만든다면, 은하에 필요한 씨앗까지 왜 지워 버리지 않을까? 답은 양자역학이다. 살아남은 요동은 다림질이 완벽하지 않아 남은 흔적이 아니다. 새로 생성되는 것이다.

잠시 팽창을 잊어 보자. 장이 거의 완벽하게 균질하더라도 양자역학은 \(\phi\)와 \(\dot\phi\)를 둘 다 정확히 고정된 값으로 두도록 허용하지 않는다. 바닥 상태에도 요동이 있다. 보통의 조화 진동자에 영점 에너지가 있는 것과 같은 이유다. 일상적인 상황에서 이런 요동을 알아보기 어려운 것은 파장이 짧고 빠르게 진동하여 관측값을 평균 내면 사라지기 때문이다.

팽창은 이야기를 바꾸어 놓는다. 공변 좌표에서 장의 모드는 좌표 파장이 대략 그대로 유지되는 반면, 물리적 파장은 스케일 인자에 따라 커진다.
\begin{equation}
\lambda_{\rm phys}=a(t)\,\Delta x.
\end{equation}
인플레이션 동안 물리적 허블 규모는 대략 일정하다.
\begin{equation}
d_H\sim H^{-1},
\end{equation}
따라서 이에 대응하는 공변 지평선은 다음과 같이 작아진다.
\begin{equation}
d_{H,\mathrm{coord}}\sim \frac{1}{aH}.
\end{equation}

밧줄 비유를 생각하면 요동이 얼어붙는 과정을 직관적으로 이해할 수 있다. 밧줄의 각 부분이 이웃한 부분에 끌리기 때문에 파동이 진동한다. 하지만 파장이 지평선만큼 커지면 관련된 이웃 부분은 인과적으로 접촉할 수 없는 곳에 있게 된다. 복원력이 사라진다. 모드가 지평선보다 길어지면 얼어붙는다.

이 과정은 계속 반복된다. 파장이 긴 모드가 먼저 얼어붙고, 짧은 모드는 나중에 얼어붙는다. 아주 작은 파장으로 시작한 모드도 결국 늘어나 지평선을 가로지른다. 따라서 인플레이션은 구조를 단순히 지워 버리지 않는다. 얼어붙은 장의 변화를 계속 만들어 낸다.

인플레이션이 오래 지속된다고 해서 이 과정이 멈추는 것도 아니다. 얼어붙은 오래된 모드는 엄청난 규모로 늘어나지만, 새로운 모드도 계속 얼어붙는다. 누적된 양상은 우주론에서 의미 있는 관점으로 통계적 정상 상태를 이룬다. 이것이 근사적으로 규모 불변인 스펙트럼의 기원이다. 같은 지평선 교차 과정이 연속된 로그 규모에서 반복되어 파장 대역이 달라져도 대략 같은 진폭을 만들어 낸다.

장이 평탄부를 벗어나면 지평선은 다시 커진다. 한때 일시적인 인플레이션 지평선 바깥에 있던 모드들이 나중에 다시 안으로 들어온다. 흔히 쓰는 표현으로는 이 모드들이 인플레이션 동안 지평선을 나갔다가 그 뒤에 다시 들어온다. 그때쯤에는 거시적 규모까지 늘어나 있으므로, 다시 진동하기 시작할 때 매우 큰 구조를 이룬다.

\begin{classroomqa}
\lecturetimestamp{01:12:02}
\audiencequestion{시공간 도표에서 지평선은 어떻게 보이나요?}
\lecturerresponse{지수 팽창 동안 지평선까지의 물리적 거리는 대략 \(c/H\)로 일정하지만, 공변 도표에서는 좌표 척도가 늘어나므로 이 일정한 물리적 거리가 점점 작아 보인다. 공변 좌표로 나타낸 지평선은 \(1/(aH)\)처럼 줄어든다.}
\end{classroomqa}

\begin{classroomqa}
\lecturetimestamp{01:15:59}
\audiencequestion{요동이 늘어날 때 진폭은 어떻게 되며, 정상 상태가 있나요?}
\lecturerresponse{개별 모드는 늘어나면서 희석되지만, 새로운 모드가 계속 얼어붙는다. 그 결과 요동이 완전히 사라지는 것이 아니라 근사적으로 규모 불변인 통계적 정상 양상이 형성된다.}
\end{classroomqa}

\section{얼어붙은 장의 변화에서 밀도 변화로}

장이 장소마다 조금씩 달라지면, 서로 다른 영역이 약간 다른 시점에 가장자리를 넘어갈 때 그 차이가 밀도 차이로 바뀐다. 조금 더 아래쪽에 있던 영역이 먼저 절벽에 도달한다. 그곳의 진공 에너지가 복사로 전환되면, 복사는 팽창에 따라 희석된다. 이웃 영역은 여전히 평탄부에 있어 진공 에너지가 지배적이며 희석되지 않는다. 두 번째 영역이 떨어질 때쯤이면 첫 번째 영역의 에너지는 이미 적색 편이를 거쳐 일부 줄어든 뒤다.

그래서 굴림이 느릴수록 최종 밀도 대비가 커진다. 장이 평탄부를 따라 아주 천천히 움직이면, 장의 작은 값 차이가 두 영역이 절벽에 도달하는 시점 사이의 더 큰 시간 차이에 해당한다. 먼저 떨어진 영역의 에너지는 나중에 떨어진 영역이 에너지를 전환하기 전까지 더 오래 희석된다. 기울기가 완만할수록 종단 속도는 작아지고 최종 밀도 변화는 커진다. 평탄부가 더 높아도 요동이 커지는 경향이 있다. 평탄부의 폭은 주로 인플레이션이 얼마나 오래 지속되는지를 좌우한다.

요동 자료가 엄격하게 제한하는 것은 높이와 기울기의 조합뿐이다. 에너지 규모가 충분히 높다면 중력파 기여분도 관측할 수 있어야 한다. 수치 한계는 대략적인 범위로만 제시한다. 대략 말해 평탄부의 에너지 밀도가 플랑크 단위로 \(10^{-14}\)보다 훨씬 높다면 우주 마이크로파 배경에서 중력파 구조가 이미 지나치게 크게 보여야 한다. 반대로 \(10^{-16}\)보다 훨씬 낮다면 그런 신호를 관측할 현실적인 가능성이 거의 없다. \editorialnote{강연자는 이 수치들이 근사값이라고 명시하고 정확한 값을 기억하지 못한다고 말한다. 여기서는 크기 차수에 관한 진술로만 남겨 둔다.}

단위는 통상적인 자연 단위계 방식으로 설명된다. 질량이 \(m\)인 입자의 콤프턴 파장은 다음과 같다.
\begin{equation}
\lambda_C=\frac{\hbar}{mc}.
\end{equation}
\(c=\hbar=1\)이면 질량과 역길이는 같은 단위를 갖는다. 에너지 밀도는 부피당 에너지이므로 질량의 4제곱 단위를 갖는다.
\begin{equation}
[\rho]=(\mathrm{mass})^4.
\end{equation}

한 지점에서 보면 얼어붙은 모드가 누적되는 과정은 무작위 보행처럼 보인다. 공간의 각 지점은 서로 다른 무작위 보행을 겪으며, 바로 이 차이가 양자 잡음을 공간적 변화로 바꾼다.

\begin{classroomqa}
\lecturetimestamp{01:33:04}
\audiencequestion{\(\phi\)는 무작위 보행을 하며, 그렇다면 더 이른 시작점은 중요하지 않게 되나요?}
\lecturerresponse{어느 한 고정된 지점에서 얼어붙은 모드들이 무작위 위상으로 차례로 더해지므로 \(\phi\)는 실제로 무작위 보행을 한다. 그런데 바로 그 성공이 어려움이기도 하다. 인플레이션이 나중의 요동으로 이전 정보를 계속 대체하므로, 인플레이션 이전 단계는 복원하기 어려워진다.}
\end{classroomqa}

\section{인플레이션이 감추는 것}

이 이론은 매우 구체적인 의미에서 성공적이다. 하나의 근본 메커니즘으로 엄청나게 넓은 규모 범위에 걸친 요동들을 서로 연관 짓는다. 그러나 같은 메커니즘이 이전에 있었던 일을 씻어 없앤다. 이전의 불규칙성은 늘어나 사라지고, 나중에 생성된 양자 요동이 중요하게 남는다. 그래서 인플레이션은 매력적이면서도 답답하다. 우주의 매끄러움과 요동 스펙트럼을 설명하지만, 초기 역사의 상당 부분도 지워 버린다.

이 때문에 시작점에 대한 불편함이 생긴다. 매우 높은 에너지 밀도가 문제인 것은 아니다. 평탄부의 에너지가 자연스럽게 플랑크 규모라고도, 아주 작다고도 볼 수 없다는 점이 문제다. 일상적인 단위로는 크지만 플랑크 단위로는 여전히 작다. 근본 이론의 명백한 첫 출발점처럼 보이지 않는다.

이 단순한 그림에는 문자 그대로의 \(t=0\)도 없다. 인플레이션이 너무 오래 지속되면 그보다 앞선 단계에 대한 정보는 거의 모두 사라진다. e-접힘 수가 필요한 최소치보다 조금만 더 컸다면 이전 단계의 흔적이 일부 남을 수도 있지만, 일반적으로 인플레이션은 과거를 지워 버린다.

드 시터 공간을 다룰 때도 같은 주의가 필요하다. 순수한 지수 법칙은 후기의 인플레이션 가지에 대해 맞지만, 정확한 드 시터 해는 다음처럼 나타내는 편이 낫다.
\begin{equation}
a(t)\propto \cosh(Ht).
\end{equation}
이 해는 최소 크기까지 수축한 다음 팽창한다. 수축하는 가지를 우주론에서 문자 그대로 받아들여야 하는지는 열린 문제로 남겨 둔다.

\begin{classroomqa}
\lecturetimestamp{01:37:31}
\audiencequestion{지수적 거동이 과거로 무한히 이어져 시작이 없었을 수도 있나요?}
\lecturerresponse{이 단순한 그림에서는 그렇지 않다. 정확한 드 시터 공간은 모든 시간에 걸쳐 영원히 팽창하는 지수적 우주가 아니다. 수축하는 가지와 팽창하는 가지를 지닌 \(\cosh(Ht)\)에 더 가깝다. 그 이전 가지를 문자 그대로 받아들여야 하는지는 열린 질문으로 남겨 둔다.}
\end{classroomqa}

\section{미세 조정, 곡률, 인간원리의 한계}

일차 근사에서 중요한 매개변수는 평탄부의 높이와 기울기, 폭 \(\Delta\phi\)이다. 기울기는 종단 속도를 정하고, 따라서 요동의 크기를 정하므로 중요하다. 폭은 인플레이션이 얼마나 오래 지속되는지를 결정하는 데 영향을 준다. 기본적인 장의 변화량은 플랑크 질량보다 훨씬 커서는 안 되고 에너지 밀도는 플랑크 단위로 1 정도를 넘지 않아야 한다는 통념이 있다. 이런 대략적인 한계 아래에서는 충분한 인플레이션을 일으키는 조건이 매개변수 공간 전체를 채우지 않는다. 여기서 인용하는 추정에 따르면 기울기를 충분히 완만하게 만들어 대략 \(60\)에서 \(62\) 사이의 e-접힘을 얻으려면 몇 퍼센트, 즉 약 \(2\%\) 수준의 미세 조정이 필요하다.

\begin{classroomqa}
\lecturetimestamp{01:46:02}
\audiencequestion{인플레이션을 매개변수 공간으로 이해할 수 있나요? 미세 조정은 어느 정도인가요?}
\lecturerresponse{그렇다. 일차 근사에서는 평탄부의 높이와 기울기, 폭을 생각한다. 강연에서는 충분한 e-접힘을 얻으려면 몇 퍼센트 수준, 약 \(2\%\)의 미세 조정이 필요하다고 대략 추산한다.}
\end{classroomqa}

이 정도의 완만한 미세 조정은 현재의 우주 상수, 약 \(10^{-123}\)이라는 훨씬 더 날카로운 문제와 비교된다. 논의는 인간원리 문제로 넘어간다. 두 문제 모두 구조 형성에 필요한 조건을 묻는 방식으로 다시 표현할 수 있기 때문이다.

양의 우주 상수는 뉴턴적 그림에서 거리에 비례하는 보편적 반발력처럼 작용한다. 그 값이 훨씬 컸다면 오늘날 은하를 산산이 흩어 놓아야만 치명적이 되는 것은 아니다. 원시 \(10^{-5}\) 요동이 은하로 응축되는 일을 막을 만큼만 커도 충분하다. 제시된 예에 따르면 관측된 값보다 대략 두 자릿수만큼 큰 값이어도 이미 그런 일이 일어난다. 이것이 와인버그의 논지와 관련된 핵심이다. 우주 상수는 정확히 0일 필요는 없고, 구조가 형성될 수 있을 만큼 작기만 하면 된다.

음의 곡률도 같은 방향으로 작용한다. 앞선 강의의 뉴턴적 표현으로는 \(k<0\)일 때 후퇴 속도가 탈출 속도보다 크다는 뜻이다. 바깥으로 향하는 추력이 지나치게 크면 작은 원시 과밀 영역도 은하로 발달하지 못한다. 인플레이션은 그 곡률을 희석한다. 인용된 하한은 약 \(59.5\) e-접힘이다. \(62\) 정도를 문턱보다 충분히 높은 대표적인 인플레이션 역사로 본다.

이렇게 하면 확률에 관한 질문이 달라진다. 충분한 인플레이션을 일으키는 매개변수 공간의 원래 부피를 묻는다면 답은 몇 퍼센트에 불과할 수 있다. 은하가 형성되고 관측자가 존재하여 질문을 던질 수 있는 우주에 있다고 조건을 붙이면 조건부 확률은 훨씬 커질 수 있다. 논의는 확정된 해답이 아니라 미세 조정 문제와 그에 대한 인간원리적 재해석, 그리고 큰 논쟁이 여전히 이어지고 있다는 인식으로 끝난다.
